\subsection{1.9.6 Planificación del
proyecto}\label{planificaciuxf3n-del-proyecto}

1.9.6.1 Resumen de la planificación del proyecto

{\def\LTcaptype{none} % do not increment counter
\begin{longtable}[]{@{}
  >{\raggedright\arraybackslash}p{(\linewidth - 8\tabcolsep) * \real{0.0689}}
  >{\raggedright\arraybackslash}p{(\linewidth - 8\tabcolsep) * \real{0.0318}}
  >{\raggedright\arraybackslash}p{(\linewidth - 8\tabcolsep) * \real{0.2767}}
  >{\raggedright\arraybackslash}p{(\linewidth - 8\tabcolsep) * \real{0.1742}}
  >{\raggedright\arraybackslash}p{(\linewidth - 8\tabcolsep) * \real{0.4484}}@{}}
\toprule\noalign{}
\begin{minipage}[b]{\linewidth}\raggedright
\textbf{ID}
\end{minipage} &
\multicolumn{2}{>{\raggedright\arraybackslash}p{(\linewidth - 8\tabcolsep) * \real{0.3085} + 2\tabcolsep}}{%
\begin{minipage}[b]{\linewidth}\raggedright
\textbf{Fase / Módulo}
\end{minipage}} & \begin{minipage}[b]{\linewidth}\raggedright
\textbf{Periodo}
\end{minipage} & \begin{minipage}[b]{\linewidth}\raggedright
\textbf{Actividades principales}
\end{minipage} \\
\midrule\noalign{}
\endhead
\bottomrule\noalign{}
\endlastfoot
\textbf{M1} &
\multicolumn{2}{>{\raggedright\arraybackslash}p{(\linewidth - 8\tabcolsep) * \real{0.3085} + 2\tabcolsep}}{%
\textbf{Definición (M1)}} & \textbf{18-feb - 1-mar} & Tema, objetivos,
metodología, planificación \\
\textbf{M2} &
\multicolumn{2}{>{\raggedright\arraybackslash}p{(\linewidth - 8\tabcolsep) * \real{0.3085} + 2\tabcolsep}}{%
\textbf{Estado del arte (M2)}} & \textbf{2-mar - 22-mar} & Revisión
bibliográfica, redacción capítulo 2 \\
\textbf{M3} &
\multicolumn{2}{>{\raggedright\arraybackslash}p{(\linewidth - 8\tabcolsep) * \real{0.3085} + 2\tabcolsep}}{%
\textbf{Implementación (M3)}} & \textbf{23-mar - 10-may} & Desarrollo
técnico completo \\
M3-1 & ~ & Preparación de datos & 23-mar - 27-mar & Limpieza inicial y
revisión de datos \\
M3-2 & ~ & Transformación supervisada & 28-mar - 4-abr & Generación de
lags, rolling windows, codificación categórica \\
M3-3 & ~ & Modelado Árboles (v1) & 5-abr - 15-abr & Entrenamiento RF,
XGBoost, LightGBM + validación temporal \\
M3-4 & ~ & Reconciliación jerárquica & 16-abr - 22-abr & Implementación
Bottom-Up y MinT \\
M3-5 & ~ & Interpretabilidad (SHAP) & 23-abr - 30-abr & Cálculo
global/local, análisis estabilidad \\
M3-6 & ~ & LSTM multiserie & 1-may - 7-may & Implementación y
entrenamiento LSTM + generación de predicciones base para posterior
reconciliación jerárquica. \\
M3-7 & ~ & Evaluación comparativa & 8-may - 10-may & Comparación de
métricas RMSE/MAE/WRMSSE + coste computacional \\
\textbf{M4} &
\multicolumn{2}{>{\raggedright\arraybackslash}p{(\linewidth - 8\tabcolsep) * \real{0.3085} + 2\tabcolsep}}{%
\textbf{Redacción de la documentación del TFM (M4)}} & \textbf{11-may -
2-jun} & ~ \\
M4-1 & ~ & Memoria preliminar & 11-may - 15-may & Redacción inicial con
resultados \\
M4-2 & ~ & Memoria final & 16-may - 24-may & Revisión final y ajustes \\
M4-3 & ~ & Presentación & 25-may - 2-jun & Vídeo y diapositivas \\
\textbf{M5} &
\multicolumn{2}{>{\raggedright\arraybackslash}p{(\linewidth - 8\tabcolsep) * \real{0.3085} + 2\tabcolsep}}{%
\textbf{Defensa del proyecto}} & \textbf{3-jun - 26-jun} & ~ \\
M5-1 & ~ & Entrega tribunal & \textbf{5-jun} & Documentación final \\
M5-2 & ~ & Defensa & \textbf{26-jun} & Exposición pública \\
\end{longtable}
}

Tabla 1. Resumen de la planificación del proyecto

1.9.6.2 Calendario de la planificación del proyecto

Figura 1. Calendario de la planificación del proyecto (formato Gantt).

La planificación permite solapamientos controlados entre fases,
iniciando tareas como interpretabilidad o reconciliación jerárquica una
vez establecida una primera versión del modelo base. Esto optimiza el
tiempo disponible y reduce riesgos de acumulación de tareas en fases
finales.
